%% frag_discussion.tex
%% Replaces the opening of Section 5 (from \section{Discussion} through the
%% end of its introductory \begin{itemize}...\end{itemize}), up to but not
%% including \subsection{Implications for Software Update Question Answering}.

\section{Discussion}\label{sec:discussion}

\subsection{What Generalizes}\label{sec:generalizes}

The finding most likely to transfer to other systems is also the simplest to
state: \emph{a rewriting agent must add to the retrieval it performs, not
substitute for it.} The distinction is easy to lose because it is not visible
in the agent's own output. A rewritten query inspected in isolation looks
correct --- it names the vendor, uses release-note vocabulary, and reads like a
well-formed search. What is invisible at that level is that issuing it
\emph{instead of} the user's question silently narrows the candidate set, and
that the narrowing is unrecoverable downstream. Our pipeline exhibited this for
its entire development history without any component appearing to malfunction.

Three properties made it hard to see, and each is a property of agentic
pipelines generally rather than of this system:

\begin{itemize}
  \item \textbf{The failure is in the wiring, not in a component.} Every agent
  did its own job correctly. The rewriter produced good rewrites; the retriever
  faithfully queried what it was given; the reranker correctly ordered what it
  received. Component-level testing cannot detect this, because no component is
  wrong.

  \item \textbf{The aggregate metric moved the right way.} The bespoke
  retrieval-quality score rose when the architecture was introduced, which is
  what a designer checks. A score computed over the retrieved set cannot report
  documents that were never retrieved --- their absence is not an observation it
  is capable of making.

  \item \textbf{A plausible cause absorbed the evidence.} Query rewriting is
  motivated in the literature by vocabulary mismatch, so a retrieval deficit
  after rewriting invites a ranking explanation. We pursued that explanation
  first, built three rerankers, and improved MRR by 150\% without materially
  moving nDCG. Only the residual gap indicated that the diagnosis was wrong.
\end{itemize}

The last point suggests a diagnostic worth generalizing. When an intervention
improves the metric it targets while leaving the outcome metric unchanged, the
intervention is likely addressing a real but secondary defect. In retrieval
specifically, MRR improving while nDCG@$k$ and Recall@$k$ stay flat is the
signature of a candidate set that does not contain the answer, and it is
distinguishable from a ranking problem only by inspecting the candidate set
directly. We therefore instrumented the pipeline to expose its pre-rerank pool,
so that ``the reranker mis-ordered it'' and ``it was never retrieved'' are
separable diagnoses rather than a single unexplained score.

\subsection{What Multi-Agent Decomposition Did and Did Not Buy}
\label{sec:decomposition}

With both defects repaired and the synthesis model held constant, the
multi-agent architecture matches the single-agent baseline on retrieval quality
and on faithfulness, at approximately twice the latency
(Section~\ref{sec:scaled-eval}). We report this as the result rather than
searching for a configuration that reverses it.

This does not establish that decomposition is worthless, and we are careful not
to claim it does. It establishes that on this benchmark, for this domain, with
these metrics, decomposition did not improve answer quality over a
well-configured single-agent baseline --- and that the improvement originally
attributed to decomposition was an artifact of the metric used to measure it.
Three qualifications bound that claim:

\begin{itemize}
  \item The baseline is not weak. It queries the same live endpoints, uses the
  same synthesis prompt, and benefits from the same reranking stage. A weaker
  baseline --- one without vendor-aware retrieval, or restricted to a single
  source --- would likely lose, but the comparison would not isolate
  architecture.

  \item Retrieval quality is not the only axis on which decomposition might pay.
  The architecture's failure traceability is real and is not captured by nDCG:
  when an answer is wrong, the pipeline localizes the fault to vendor
  extraction, query reformulation, retrieval coverage, evaluation, or synthesis.
  For a system intended to be maintained rather than merely benchmarked, this
  has value that our metrics do not score.

  \item The evaluator agent's ability to decline is likewise unscored here.
  Returning ``the retrieved sources do not answer this'' is preferable to a
  confident wrong version number in a maintenance context, and the CRAG-style
  scoring we implement rewards exactly that asymmetry --- but our ground-truth
  coverage is too small for it to carry weight.
\end{itemize}

\subsection{Metric Design as a Threat to Architectural Claims}
\label{sec:metric-threat}

The 17.2\% improvement was not fabricated; it is what the score we defined
reports, and it reproduces. The problem is what the score is capable of
observing. A measure of overlap between a question and the documents retrieved
for it rewards retrieving topically adjacent material, and a rewriting agent
that expands a question into vendor and release-note vocabulary produces
precisely that. The metric and the mechanism were, in effect, correlated
through the same vocabulary.

We think this is a general hazard for agentic-system evaluation rather than a
mistake peculiar to this work. Domain-specific systems frequently define
domain-specific quality measures, with good justification: standard IR metrics
require relevance judgments, which are expensive. But a bespoke measure designed
by the same people who designed the architecture will tend to reward the
behaviors the architecture produces. The practical mitigation is not to abandon
domain measures but to report at least one metric whose definition is
independent of the system --- pooled-judgment IR metrics, or a deterministic
benchmark score --- and to treat a disagreement between them as a finding
requiring explanation rather than a discrepancy to be reconciled toward the
preferred result.

Our answer-quality measurements carry the same lesson in a second form.
Comparing the multi-agent system's structured template against the baseline's
prose showed a 0.23 faithfulness deficit that disappeared entirely when the same
retrieved documents were synthesized through the baseline's own prompt
(Section~\ref{sec:scaled-eval}). An LLM judge comparing structured output
against fluent prose is partly scoring fluency. Any evaluation of an agentic
system that emits structured output against a baseline that emits prose is
exposed to this, and detecting it required constructing an arm for no purpose
other than removing the confound.

\subsection{Observations That Survive}\label{sec:survives}

Several observations from the original evaluation are unaffected by the
corrections above, because they do not depend on the architecture comparison.

\begin{itemize}
  \item \textbf{Heterogeneous sources provide complementary evidence.} Release
  notes give authoritative version and fix information; advisories give
  vulnerability details and mitigations; community discussions give deployment
  experience, regressions, and workarounds that have not reached official
  documentation. Users frequently ask about the operational effects of an update
  rather than its documented contents: a release note may confirm a defect was
  fixed while community discussion reveals users still experiencing it.
  Treating these as complementary rather than competing, and preserving the
  distinction in the retrieved set, improves coverage without conflating
  provenance.

  \item \textbf{Source grounding constrains version-specific claims.} Across the
  version-specific evaluation, generated answers did not contain version
  numbers, release dates, or patch identifiers absent from the retrieved
  documents. Incorrect version information has direct operational consequences
  for update planning and security decisions, so this property matters
  independently of retrieval ranking. The sample is small and this is an
  observation rather than a guarantee.

  \item \textbf{Retrieval-level adaptation operates without retraining.} Query
  expansion terms accumulate weights from retrieval outcomes and bias subsequent
  reformulation, with no human preference labels and no parameter updates. The
  learned terms are interpretable on inspection. We are more cautious than the
  conference version about the size of the effect: the tertile comparison that
  produced the 9.3\% figure is confounded with the live corpus drifting over the
  course of a run, and we did not control for it. The mechanism is sound; the
  magnitude is not established.
\end{itemize}
